\documentclass{article}
\usepackage[utf8]{inputenc}
\usepackage{geometry}
\geometry{
	a4paper,
	total={170mm,257mm},
	left=20mm,
	top=20mm,
}
\usepackage{graphicx}
\usepackage{titling}
\usepackage[american,RPvoltages]{circuiTikz}
\usepackage{tikz, tikz-cd}
\usepackage{pgfplots}
\usepackage{siunitx}
\usepackage{float}
\usepackage{amsmath,amsfonts,stmaryrd,amssymb} % Math packages
\title{The Analysis and Design of Linear Circuits  }
\author{Rodrigo Castillejos}
\date{July 2024}

\usepackage{fancyhdr}
\fancypagestyle{plain}{%  the preset of fancyhdr 
	\fancyhf{} % clear all header and footer fields
	\fancyfoot[L]{\thedate}
	\fancyhead[L]{Problem Resolution}
	\fancyhead[R]{\theauthor}
}
\makeatletter
\def\@maketitle{%
	\newpage
	\null
	\vskip 1em%
	\begin{center}%
		\let \footnote \thanks
		{\LARGE \@title \par}%
		\vskip 1em%
		%{\large \@date}%
	\end{center}%
	\par
	\vskip 1em}
\makeatother

\usepackage{lipsum}  
\usepackage{cmbright}

\begin{document}
	
	\maketitle
	
	\noindent\begin{tabular}{@{}ll}
		Author & \theauthor\\
		Chapter &  Basic Circuit Analysis\\
		Problem & 2.64
	\end{tabular}
	
	\section*{Exercise 2-64}
	You wish to drive a \qty{1}{\kohm} load from your car battery as shown in figure 1 . The load needs \qty{5}{\volt} across it to operate correctly . Where should the wiper on the potentiometer be set ($R_x$) to obtain the desired output voltage?
	
	\begin{figure}[H]
		\centering	
		\begin{tikzpicture}[scale=0.8]
			\ctikzset{resistors/width=1.5, resistors/zigs=9}
			\draw (0,0)   to[V, l={\qty{12}{\volt}}] (0,5)
				  to[short] ++(4,0)
				  to[potentiometer, name=P, -*] ++(0,-5);
		    \draw (P.wiper) node[xshift=-40]{$\qty{5}{\kohm}$};
			\ctikzset{resistors/width=0.8, resistors/zigs=3}
			\draw (P.wiper) to [short] ++(3,0)
				  to[R, v={\qty{5}{\volt}}, l={\qty{1}{\kohm}}] ++(0,-2.5)
				  to[short] (0,0);
			\draw[|<->, 
			transform canvas={yshift=-1.5mm, xshift=1.5mm}] (4.7,0.4) -- node[fill=white, inner sep=0pt] {$R_{\mathrm{x}}$} (P.wiper);		 
		\end{tikzpicture}
		\caption{Amplificador cascode}	
		\label{fig:figura1}
	\end{figure}
	
	\section*{Answer}
	
	Replacing the circuit of figure 1 with the model for the potentiometer that consists of two resistors we get
	\begin{figure}[H]
	\centering	
	\begin{tikzpicture}[scale=0.8]
		\draw (0,0) to[V, l={\qty{12}{\volt}}] (0,6)
			  to[short] ++(5,0)
			  to[R, -*, l={$\alpha R_p$}, i>_=$I_1$] ++(0,-3) coordinate(p1)
			  to[R, l={$(1-\alpha) R_p$}, -*, i>_=$I_2$] ++(0,-3);
		\draw (p1) to[short] ++(4,0)
			  to[R, l={\qty{1}{\kohm}}, v={\qty{5}{\volt}}, i>^=$I_3$] ++(0,-3)
			  to [short] (0,0);
		\draw (p1) node[above right]{$V_{1}$};
	\end{tikzpicture}
	\caption{Amplificador cascode}	
	\label{fig:figura2}
\end{figure}	
	
	Applying the KCL on node $V_1$ we get
	
	\begin{align}
		I_1 = I_2 + I_3 \nonumber
	\end{align}
	
	\begin{align}
		\frac{\qty{12}{\volt} - \qty{5}{\volt}}{\alpha R_p} = \frac{\qty{5}{\volt}}{(1-\alpha)R_p} + \frac{\qty{5}{\volt}}{\qty{1}{\kohm}} \nonumber \\
		\frac{\qty{7}{\volt}}{\alpha R_p} = \frac{\qty{5}{\volt}}{(1-\alpha)R_p} + \qty{5}{\milli\ampere} \nonumber \\
		\frac{\qty{7}{\volt}}{\alpha R_p} - \frac{\qty{5}{\volt}}{(1-\alpha)R_p} = \qty{5}{\milli\ampere} \nonumber
	\end{align}
	
	\begin{align}
		\frac{7(1 - \alpha) - 5\alpha}{\alpha (1 - \alpha) R_p} = \qty{5}{\milli\ampere} \nonumber \\
		\frac{7 - 7\alpha - 5\alpha}{\alpha (1 - \alpha) R_p} = \qty{5}{\milli\ampere} \nonumber \\
		\frac{7 - 12 \alpha}{\alpha (1 - \alpha) R_p} = \qty{5}{\milli\ampere} \nonumber \\
		7 - 12\alpha = 0.005\alpha (1-\alpha)R_p \nonumber \\
		7 - 12\alpha = 25\alpha - 25 \alpha ^2 \nonumber \\
		25 \alpha ^2 - 37\alpha + 7 = 0
	\end{align}
	
	Solving the quadratic function yields
	
	\begin{align}
		\alpha_1 = 1.2573 \nonumber \\
		\alpha_2 = 0.2227 \nonumber
	\end{align}
	
	We know the alpha value cannot be values greater than 1, so the solution for this equation is $\alpha=0.2227$. Therefore, the value of $R_x$ must be
	
	\begin{align}
		R_x = (1 - \alpha) R_p = (1 - 0.2227)(\qty{5}{\kohm}) = \qty{3886.5}{\ohm}
	\end{align}
	
\end{document}